\begin{abstract}
人工智能社区一直在通过开发“语言智能体”来探索通向通用人工智能（AGI）的路径。语言智能体是同时涉及提示技术与工具使用方法的复杂大语言模型（LLM）流水线。尽管语言智能体已在许多现实任务上展现出令人印象深刻的能力，当前语言智能体研究仍存在一个根本局限：它以模型为中心，或者说以工程为中心。也就是说，语言智能体在提示、工具和流水线方面的进步，需要人类专家投入大量手工工程工作，而不是从数据中自动学习。我们认为，从以模型或工程为中心转向以数据为中心，也就是让语言智能体具备在环境中自主学习和演化的能力，是它们有可能实现 AGI 的关键。

本文提出\textit{\baby}：一个使用\textit{符号优化器}、让语言智能体以数据为中心自主优化自身的系统框架。具体而言，我们把智能体视为符号网络，其中可学习的权重由提示、工具以及它们的堆叠方式定义。智能体符号学习通过模拟联结主义学习中的两种基础算法——反向传播与梯度下降——来优化语言智能体内部的符号网络。与数值权重不同，\baby{}处理的是权重、损失和梯度在自然语言中的对应物。我们在标准基准和复杂现实任务上开展概念验证实验，结果表明，\baby{}使语言智能体能够在创建并部署到真实环境之后继续更新自身，从而形成“自演化智能体”。我们展示了\baby{}框架的潜力，并将整个框架开源，以促进未来对\textit{以数据为中心}的智能体学习研究。
\end{abstract}
